\section{Projection of the primitive systems}\label{sec:rg-cover-comparison}

We construct the primitive systems on one large torus and
project them to smaller tori through odd covering maps. Projection preserves the
estimates with a fixed loss in the common constant and
one loss in polymer decay. The comparison of local steps
then uses supported rows near a common lift, where the
signed image sum has only one surviving term.
This is the construction of \cite[\S\S15.1--15.3]{QED3Master}.

\subsection{Odd covers and projected kernels}
\label{subsec:rgcov-projection}

Fix $k\ge0$, $0<\mu\le1$, $0<e_k\le\bar e_0$, and admissible
sides $\ell_2=J_c\ell_1$, where $J_c$ is odd.
Write $T_i=\eta_k\Z^3/\ell_i\Z^3$ and
$p:T_2\to T_1$ for the cover. Let
$\operatorname{pr}_i$ denote the covering-lattice projections.
Pullback $p^*$ is inclusion of functions on the covering
index sets. They have character
$\chi(w)=1$ for photon fields, coordinates and parameters,
and $\chi(w)=(-1)^{w_1+w_2+w_3}$ for spinors and pair
fields. Oddness gives
$\chi(J_cw)=\chi(w)$, so a field on $T_1$ is a field
on $T_2$ with the same character.
The image is
\[
 \mathrm{Per}=\{f:f(t+\ell_1w)=\chi(w)f(t),\ w\in\Z^3\}.
\]
Translations move the photon site and preserve the chamber
label. In particular pullback commutes with $J,J^*$,
complex conjugation and constant spin matrices.

The index weight is $1$ on unit sites, unit bonds and
photon coordinate indices, $\eta_k^3$ on fine photon
sites and bonds, and $\eta_k^3/48$ on chamber sites.
For fields with a common index set and character,
\begin{equation}\label{eq:rgcov-pullback-pairing}
 \|p^*f\|_\infty=\|f\|_\infty,\qquad
 \langle p^*f',p^*f\rangle_{T_2}
                                      =J_c^3\langle f',f\rangle_{T_1}.
\end{equation}
Indeed a fundamental set for $T_2$ is the union of
$\mathrm{Fnd}_1+\ell_1w$, over representatives
$w\in(\Z/J_c\Z)^3$. The two characters cancel in
each paired summand, and the weight is the same on both
tori. At zero the chamber weight remains $1/48$.

Pullback maps the parameter polydisk, including its interior,
into the corresponding larger polydisk, since
$(p^*s)_{\mathrm{pc}}=s_{p\,\mathrm{pc}}$.
It also maps
$\widetilde{\mathcal G}_k(T_1)$ into
$\widetilde{\mathcal G}_k(T_2)$.
To check the latter assertion, write the domain decomposition
$A=A_0+A_1$. Pullback preserves reality of $A_0$ and
the sup norm of $A_1$. Each five-cube neighborhood has
a common nonwrapping lift of side $11M<\ell_1$.
The real local gauge representative of $A_0$ on its
image in $T_1$ pulls back on this lift to the
corresponding neighborhood in $T_2$, with its same
sup bound and gradient identity. This is precisely
the local real-gauge condition in the inverse domain.

Every exact finite fermion operator used below intertwines
with pullback. For the chamber Dirac operator this follows
term by term from its four neighbors and bond phases.
For $\mathcal Q_k=Q_kJ^*$ and its reverse it follows from
the same block sums and the same six paths on the cover;
$J^*$ averages the same $48$ labels.
The penalty and the finite matrices multiplying the
resolvent have the same finite block formulas.
Sums and products preserve intertwining.
If $O_2p^*=p^*O_1$ and both operators are invertible,
then
$O_2p^*O_1^{-1}=p^*$ and hence
$O_2^{-1}p^*=p^*O_1^{-1}$.
The domain preservation just proved makes the already
established inverse bounds applicable on both tori.
It follows, from the formulas in Section~\extref{4}{FermionBlockManuscript},
that the fine resolvents, $H_k,\bar H_k,D_k,M_k,\Gamma_k$
intertwine at the same connection:
\begin{equation}\label{eq:rgcov-exact-fermion-intertwining}
 O_2(p^*A)p^*=p^*O_1(A).
\end{equation}
At zero use the entire chamber Dirac stencil and the
inverse of $F_y$, on all chamber components.
The exact photon operators have the corresponding
identity \exteq{46}{MaxwellBlockManuscript}.

Let a primitive system on $T_2$ be given.
For a parameter family $O_2(s,A)$ put
$\widetilde O(s,A)=O_2(p^*s,p^*A)$.
For a polymer family put
\[
 \widetilde O(X,A)=\sum_{X'\in\mathcal P_k(T_2):pX'=X}
                                                    O_2(X',p^*A).
\]
This is applied to $M_k$ for $k\ge1$ and to $S_{k,y}$
at every level. A projected polymer is connected:
adjacent $M$-cubes have equal or adjacent images.
Moreover
\begin{equation}\label{eq:rgcov-projected-polymer-length}
                           d_M(pX')\le d_M(X').
\end{equation}
A projected net meets all image cubes. Projecting its chosen
finite segment representation preserves the lengths of the segments;
taking the infimum gives \eqref{eq:rgcov-projected-polymer-length}.
There are only finitely many polymers on either torus.

\begin{proposition}\label{prop:rgcov-kernel-projection}
If the system has translation covariance under
$\ell_1\Z^3$, its projection is well defined by
$p^*O_1=\widetilde Op^*$ and has kernel
\begin{equation}\label{eq:rgcov-signed-image-kernel}
 O_1(\cdot)(t,t')
  =\sum_{w\in(\Z/J_c\Z)^3}\chi(w)
                       \widetilde O(\cdot)(t,t'+\ell_1w).
\end{equation}
This formula uses the same index weights on both tori.
The signed summand depends only on the class of $w$.
Transposition of an intertwining operator also intertwines:
\begin{equation}\label{eq:rgcov-transpose-intertwining}
                          O_2^Tp^*=p^*O_1^T.
\end{equation}
\end{proposition}

\begin{proof}
The pulled-back parameters and connection are invariant
under $\ell_1$ translations. Parameter-family covariance
therefore says that $\widetilde O$ commutes with those
translations. For polymer families, translation permutes
the polymers $X'$ with $pX'=X$, so the same commutation
holds after their finite sum.
An operator commuting with these translations maps the
character eigenspace $\mathrm{Per}$ into itself.
Since $p^*$ is injective onto that space, this defines
exactly one projected operator.

Sum its action over the fundamental sets used in
\eqref{eq:rgcov-pullback-pairing} and insert
$f(t'+\ell_1w)=\chi(w)f(t')$.
The coefficient of $\nu(t')f(t')$ is
\eqref{eq:rgcov-signed-image-kernel}.
Replacing $w$ by $w+J_cv$ multiplies the character
by $\chi(v)$ and the large-torus kernel by the same
character, so the product is unchanged.
This proves the asserted kernel formula, including
the weights $\eta_k^3/48$ or $1/48$ on chamber columns.

Translations have transpose $T_v^T=T_{-v}$ for the
displayed bilinear pairings. Hence $O_2^T$ also commutes
with $\ell_1$ translations and maps $\mathrm{Per}$
into itself. Write $O_2^Tp^*f'=p^*g$.
For every source field $f$, intertwining and
\eqref{eq:rgcov-pullback-pairing} give
\[
 J_c^3\langle g,f\rangle_{T_1}
 =\langle O_2^Tp^*f',p^*f\rangle_{T_2}
 =J_c^3\langle f',O_1f\rangle_{T_1}.
\]
Nondegeneracy gives $g=O_1^Tf'$.
\end{proof}

\subsection{The uniform projection estimate}
\label{subsec:rgcov-projection-estimates}

The supports of primitive row and column indices have
sup diameter at most $L$. Assume
$\gamma_wM\ge1$. Then
\begin{equation}\label{eq:rgcov-exponential-image-sum}
 \sum_{w\in(\Z/J_c\Z)^3}
     \exp\big(-\gamma_wd_2(t,t'+\ell_1w)\big)
                     \le64\exp\big(-\gamma_wd_1(t,t')\big).
\end{equation}
To prove this, translate the column lift by an
$\ell_1$ period so that the distance of the two support
sets equals $d_1\le\ell_1/2$.
For each class of $w$, choose a representative $w'$
which realizes its distance on $T_2$. Distinct classes
have distinct representatives. Their distances are
at least $d_1$; for $|w'|_\infty=j\ge2$ they are at least
\[
 \ell_1j-d_1-2L
                         \ge d_1+(\ell_1-2L)(j-1).
\]
This uses the support diameters and $d_1\le\ell_1/2$.
Since $\ell_1\ge13LM$, we have
$\gamma_w(\ell_1-2L)\ge6$.
There are $27$ lattice vectors of norm at most $1$
and at most $26j^2$ of norm $j$.
The sum is therefore bounded by
\[
 \exp(-\gamma_wd_1)
                 \left(27+\sum_{j\ge2}26j^2\exp(-6(j-1))\right).
\]
The latter series starts below $0.269$ and its
successive ratios are at most
$(9/4)\exp(-6)<1/100$. Thus the bracket is below
$27.3<64$, proving \eqref{eq:rgcov-exponential-image-sum}.

There is a second, polymer summation estimate. For a
unit or chamber row index $t$, let $\square_t$ be its
$M$-cube on $T_2$. By
\eqref{eq:rgcov-projected-polymer-length} and counting,
\begin{equation}\label{eq:rgcov-polymer-image-sum}
 \sum_{\substack{pX'=X\\X'\supseteq\square_t}}
    \exp\big(-(\kappa_w+\kappa_0)d_M(X')\big)
                 \le K_0\exp\big(-\kappa_wd_M(X)\big).
\end{equation}
Indeed the factor with rate $\kappa_w$ is at most
the factor on the right; the remaining sum is bounded
by $\sum_{X'\supseteq\square_t}\exp(-\kappa_0d_M(X'))\le K_0$.
The additional top-system decay $\kappa_0$ pays for the
polymer image sum.

\begin{proposition}\label{prop:rgcov-projected-system}
Suppose the system on $T_2$ has the primitive properties
with parameters
\[
 (\alpha_0,\gamma_w,\kappa_w+\kappa_0;C_w^{(1)}),
 \qquad C_w^{(1)}\ge4\exp(2\gamma_w)M^{27\alpha_0}.
\]
Its projection has the corresponding properties on $T_1$
with
\begin{equation}\label{eq:rgcov-common-constant}
                          C_w=64K_0C_w^{(1)}
\end{equation}
and parameters $(\alpha_0,\gamma_w,\kappa_w;C_w)$.
In particular it satisfies the primitive hypotheses of
Theorem~\extref{A}{NonlinearStep}.
At zero the resolvent estimate retains $(1+y)^{-1}$.
\end{proposition}

\begin{proof}
We verify the identities, estimates and local structure
needed for the step.
At full strength the exact intertwining identities and
the primitive sum identities on $T_2$ give
\[
 H_{k,1}^{\mathrm{ph}}(1)=H_{k,1}^{\mathrm{ph}},\quad
 C_{k,1}^{1/2}(1)=C_{k,1}^{1/2},\quad
 H_{k,1}(1,A)=H_{k,1}(A),\quad
 \Gamma_{k,1}(1,A)=\Gamma_{k,1}(A).
\]
For the two polymer families, sum first over all $X$:
each $X'$ occurs once, under $X=pX'$.
Equation~\eqref{eq:rgcov-exact-fermion-intertwining}
then identifies the sum with the exact operator on $T_1$.
All projected entries are finite sums of original entries
composed with the linear pullback. Continuity on the
closed parameter polydisk, holomorphy in its interior and
in the connection for any fixed parameter, and Borel
dependence on the resolvent parameter follow from these
finite sums.

For the photon reconstruction, its first and H\"older
derivative kernels have the same image-sum formula:
the difference acts in the row variable, and short
lifted differences agree on the two tori.
Equation~\eqref{eq:rgcov-exponential-image-sum} bounds
each by $64C_w^{(1)}\exp(-\gamma_wd_1)$.
It gives the same bound for the square-root kernel
and for $\Gamma_k$.
The $H_k$ sup and short-pair H\"older bounds instead
follow by the isometric pullback in
\eqref{eq:rgcov-pullback-pairing}; their constant is
$C_w^{(1)}$. A short pair and its six paths lift
identically. An intracluster quotient is
$\eta_k^{-\alpha}(f(c,g')-f(c,g))$ on each torus.

For $M_k(X,A)$, only $X'$ containing $\square_t$
contribute to its row, by original polymer support.
Combine \eqref{eq:rgcov-signed-image-kernel},
\eqref{eq:rgcov-exponential-image-sum} and
\eqref{eq:rgcov-polymer-image-sum} to get
\[
 |M_{k,1}(X,A)(t,t')|_c
     \le C_w\exp(-\gamma_wd_1(t,t')-\kappa_wd_M(X)).
\]
For a positive-level fine resolvent, the pullback of
a unit-block indicator is the disjoint sum of its
$J_c^3$ translated large-torus block indicators.
The same two estimates bound its block operator
norm. At zero use its chamber kernel directly:
the original factor $(1+y)^{-1}$ is common to all
terms of both image sums and remains on the right.
All column weights are unchanged, so no further
factor $48$ appears.

For component locality, an index of one parameter kind
whose support is contained in its active union belongs
to exactly one link-component. Here is the geometric
point. Site supports are singletons. For a bond of
length $a\in\{\eta_k,1\}$, cubes of radius $M$ meeting
its support either share an endpoint or would have
centers differing in that direction by $2M+a$.
The latter is impossible since the centers differ by
a multiple of $M$ and $0<a<M$.
For a coordinate index, its support has diameter at
most $L$. The centers of two cubes of radius $LM$
meeting it differ coordinatewise by less than $3LM$,
and are multiples of $LM$; choosing the coordinate
midpoint when the difference is $2LM$, and either
center when it is $LM$, exhibits a lattice site in
both cubes. Thus all cubes meeting the support are
linked. The same argument applies on either torus
by lifting the support and these cubes into a
nonwrapping neighborhood.

Pull back the active parameter cubes on $T_1$.
A link-component on $T_2$ maps into one component
on $T_1$, since the covering map preserves links.
For a row in a small-torus component, the unique
large-torus component containing its lift maps into
that component. Original component restriction
confines all surviving columns there.
Changing the parameters and connection outside the
small component leaves these original rows unchanged;
the signed image sum therefore preserves their row
dependence. The identical argument applied to each
lifted column preserves column component restriction.
Each row-vanishing property stated for an original family passes
to its projection. In particular, if a cube meeting $K_c$ is
missing, its unique lift meeting the lifted $K_c$ is missing
too, so the $H_k$ row vanishes. For $k\ge1$, the
parameter-independent term $\mathcal B_{k,0}$ in $\Gamma_k$
is retained; it preserves every union of partition $M$-cubes
\cite[Proposition~4.3 and its proof]{FermionBlockManuscript}.
For a projected polymer family a nonzero term has
both original indices in $X'$, so both projected
indices lie in $X$. Equality of connections on bonds
meeting $X$ gives equality on bonds meeting every
such $X'$, preserving the polymer's local dependence.

Gauge covariance passes through pullback because
the local gradients, restriction to unit sites and
phases commute with it. Lattice covariance follows
by lifting a symmetry $(S,a)$ with one representative
$a\in LM\Z^3$ to the same affine map on both tori;
it permutes the image polymers. A change of representative
multiplies source and target spinors by the same sign,
which cancels in the operator relation.
Transposition passes by
\eqref{eq:rgcov-transpose-intertwining};
the constant $\mathsf C$ commutes with pullback.
Complex conjugation and $\varkappa$ also commute
with it. These facts prove the charge and reality
relations for the projected primitives.

At zero, the photon and fermion diagonal factors
are unchanged by projection. Cubes containing a
given bond, and cubes meeting $K_c$, have unique
corresponding lifts, so their parameter products
are the same on both tori. The fermion factor is
the same for all $48$ labels at $c$.
Their constants are bounded by
$4\exp(2\gamma_w)M^{27\alpha_0}\le C_w^{(1)}$.
The remaining zero-level properties are those
just proved for the square root, chamber covariance
and local resolvent. This completes the system
verification with the common constant
\eqref{eq:rgcov-common-constant}.
\end{proof}

\subsection{A common top-torus system}
\label{subsec:rgcov-top-system}

On a chosen top torus use the photon families
\exteq{37}{MaxwellBlockManuscript} and the local
fermion families \exteq{57}{FermionBlockManuscript}.
In the fermion walk construction take the polymer
decay rate $\kappa'=\kappa_w+\kappa_0$.
The formulas do not depend on this rate; it specifies
the parameter restrictions and the constants in their
estimates. We write $\gamma^F=\gamma$ for the fermion
decay rate in \exteq{31}{FermionBlockManuscript}. Define
\begin{gather}
 \alpha_0=1/54,\qquad
 \gamma_w=\min(\gamma^F,\gamma^P),\quad
                                \gamma^P=\gamma_P/2, \notag\\
 C^P=116C_P^{\mathrm{ph}}M^{402\alpha_0},\qquad
 C_w^{(1)}=\max\{C^F,C^P,4\exp(2\gamma_w)M^{27\alpha_0}\},
 \qquad C_w=64K_0C_w^{(1)}.             \label{eq:rgcov-top-constants}
\end{gather}
Here $C^F,\gamma^F$ are the constants of the fermion
construction at the indicated polymer rate.
The fixed $M$ satisfies both its fermion condition and
\exteq{38}{MaxwellBlockManuscript}; the charge satisfies
the inverse and walk conditions of those constructions.

These families form a system with
$(\alpha_0,\gamma_w,\kappa_w+\kappa_0;C_w^{(1)})$.
This follows by matching their proved properties.
The photon kernel multipliers are polynomials in their
parameter cubes; at full strength the partition of unity
and the products give the exact kernels. Their path
products confine rows and columns to their active
link-components. The photon derivative and square-root
bounds have constants $C^P$ and rate $\gamma^P$.
The fermion families are the convergent weakened walks
and the localized block expressions, with their proved
support, covariance, analyticity and H\"older estimates.
Replacing their rates $\gamma^F,\gamma^P$ by the smaller
$\gamma_w$ and their constants by the larger $C_w^{(1)}$
preserves every inequality. The row and column
restriction identities contain no decay constant and
are unchanged.

At level zero the two diagonal families are used.
The photon derivative bound has constant
$4\exp(2\gamma_w)M^{27\alpha_0}$, and the fermion
diagonal factor is bounded by
$M^{8\alpha_0}$, below that constant.
There is no primitive short pair with $0<d_X<1$ at $\eta_0=1$.
The covariance and local resolvent are chamber kernels
with measure $1/48$, and the resolvent bound includes
$(1+y)^{-1}$. A unit block has total chamber measure
$1$, since
$(\eta_k^3/48)(48\eta_k^{-3})=1$.
Consequently the positive-level block estimates and
zero-level chamber estimates use precisely the common
constants \eqref{eq:rgcov-top-constants}.
Projecting this system by \cref{prop:rgcov-projected-system}
gives a system with these primitive bounds and the same $C_w$ on every
odd smaller torus. The top system itself has these
weaker bounds as well. None of these constants depends
on the level, side, mass, charge or harmonic background.

